\documentclass[11pt]{article}
\usepackage{icmlko}
\begin{document}

\title{백스물여섯 건이 스물일곱 가지 설명으로 무너진다 --- 태거 분해능은 총량으로는 병목이 아니지만 한 도메인에서는 전부다}
\author{월드모델 실험실}
\date{노트 530 · 2026년 8월 4일}
\maketitle

\begin{abstract}
노트 529 가 단일 축의 \textbf{동점률}을 재다가 ``그 축이 얼마나 거친가''를
봤다. 이걸 \emph{레코드 쌍} 수준으로 올리면 \textbf{어떤 모델도 못 넘는
천장}이 나온다 --- 특징 벡터가 같은 두 레코드는 아무도 순서를 못 매긴다.

\begin{center}
\begin{tabular}{lrrrrr}
\hline
도메인 & 유보 & 서로 다른 & 분해능 & 구별 못 쌍 & 실제 \\
\hline
\textbf{시장팝업} & $126$ & $\mathbf{27}$ & $\mathbf{0.21}$ & $\mathbf{12.7\%}$ & $54.2\%$ \\
만화 & $258$ & $186$ & $0.72$ & $0.9\%$ & $61.5\%$ \\
팝업 & $65$ & $56$ & $0.86$ & $1.2\%$ & $60.2\%$ \\
나머지 여덟 & --- & --- & $0.93\!\sim\!1.00$ & $\approx 0$ & $62\!\sim\!73\%$ \\
\hline
\end{tabular}
\end{center}

\textbf{총량으로는 반증이다} --- $100\%$ 에서 모자란 $32.8$ \%p 중 태거 몫은
$\mathbf{0.3}$ \textbf{\%p} 이고 나머지 $32.5$ \%p 가 모델 몫이다.
\textbf{그런데 한 도메인에서는 그것이 전부다.}
\end{abstract}

\begin{figure}[t]
\centering
\includegraphics[width=\textwidth]{figs/f1.pdf}
\caption{\textbf{가} --- 분해능 대 정확도. \textbf{나} --- 부족분 분해.
\textbf{다} --- 시장팝업.}
\end{figure}

\section{사전 등록 둘, 하나는 조건부 통과 하나는 반증}

\paragraph{① 분해능↔정확도 순위상관 $\ge +0.5$ --- 통과, 그러나 꼬리가 전부다.}
전체 $+0.684$ 이고 어느 한 점을 빼도 살아남는다(시장팝업 빼면 $+0.567$,
만화 빼면 $+0.580$, 팝업 빼면 $+0.580$). 그런데

\begin{quote}
\textbf{분해능 $0.9$ 이상인 여덟만 보면 $\mathbf{+0.062}$ 다.} 상관은 전부
\emph{낮은 꼬리}가 만든다.
\end{quote}

\begin{quote}
\textbf{조항} --- \textbf{분해능은 문턱 지표다.} 낮으면 그 도메인의 태거가
병목이고, 충분히 높으면 더 높다고 좋아지지 않는다. 순위를 매기는 데 쓰면 안
되고 \textbf{거르는 데} 쓴다.
\end{quote}

\paragraph{② 시장팝업과 도서가 하위 셋 --- 반증($1/2$).} 시장팝업은 맞았지만
\textbf{도서는 $0.96$ 으로 높다.} 노트 529 에서 도서의 단일 축(\texttt{grp})
동점률이 $92.9\%$ 로 제일 높았는데도 그렇다 --- \textbf{축 하나가 거친 것과
도메인이 거친 것은 다르다.} 도서는 다른 축이 많아서 벡터가 갈린다.

\section{시장팝업은 산수로 딱 맞는다}

\begin{center}
\begin{tabular}{lrr}
\hline
관측되는 축 & 서로 다른 값 & 덮음 \\
\hline
\texttt{goods\_scale} & $2$ & $94\%$ \\
\texttt{target\_breadth} & $2$ & $100\%$ \\
\texttt{venue\_prominence} & $7$ & $100\%$ \\
\hline
\end{tabular}
\end{center}

$2 \times 2 \times 7 = 28$ 이고 실측이 $27$ 이다(한 조합이 안 나타난다).
\textbf{시장팝업 태거가 채우는 축은 셋뿐이고, 그 셋이 각각 두 값 $\cdot$ 두 값
$\cdot$ 일곱 값이다.} 견줌 --- 웹툰은 관측되는 축이 $17$ 개, 만화는 $8$ 개
(그중 넷은 덮음 $3\!\sim\!21\%$)다.

\begin{quote}
그래서 시장팝업 유보 쌍의 $12.7\%$ 는 \textbf{원리상 채점이 안 된다}
--- 두 기획의 설명이 글자 그대로 같으니 어느 쪽이 나은지 물을 수가 없다.
천장이 $93.6\%$ 이고 실제가 $54.2\%$ 다.
\end{quote}

\section{그래도 병목은 여전히 모델이 아니다}

주의해서 읽어야 한다. 천장 $93.6\%$ 에 실제 $54.2\%$ 라면 \emph{시장팝업에서도}
모자란 $39.4$ \%p 중 태거 몫은 $6.4$ \%p 뿐이다. 나머지는 ``구별은 되는데 못
맞힌다''이고, 그건 축이 \emph{두 값}밖에 없어서 구별이 \emph{거의} 안 되는
것과 같은 뿌리다.

\begin{quote}
\textbf{정직하게} --- ``구별 불가능''은 태거 문제의 \emph{극단}만 센다.
``값이 두 개뿐''은 구별 불가능으로 안 잡히지만 거의 같은 병이다. 그래서
$0.3$ \%p 라는 총량은 \textbf{태거 문제의 하한}이지 전부가 아니다.
\end{quote}

\textbf{교란도 적어 둔다} --- 유보 행수와 분해능의 상관이 $+0.332$, 유보 행수와
정확도가 $+0.445$ 다. 작은 도메인이 분해능도 낮고 정확도도 낮은 성분이 섞여
있다.

\section*{그래서 무엇을 하나}

\begin{itemize}
\item \textbf{시장팝업 태거가 다음 수집 항목의 첫째다.} 지금 축 셋을 다섯로
늘리고 값 단계를 쪼개면 $27$ 그룹이 수백으로 간다. 노트 505 가 ``시장팝업
태거(칸 $15$)''를 사람이 정할 일로 적어 둔 것과 같은 자리이고, 이제
\textbf{얼마나 급한지가 숫자로 있다}.
\item \textbf{만화는 둘째다}($0.72$). 관측되는 축 여덟 중 넷이 덮음
$3\!\sim\!21\%$ 라 실질 축이 넷이다.
\item \textbf{나머지 여덟은 태깅으로 못 고친다} --- 분해능 $0.93$ 이상이라
구별은 이미 된다. 거기서는 행을 더 모으는 것이 답이다(노트 506 $\cdot$ 508).
\end{itemize}

\section*{남은 것}

\begin{itemize}
\item \textbf{``값 단계 수''를 분해능 대신 재 본다} --- 구별 불가능은 극단만
세므로, 축마다 서로 다른 값의 수를 도메인 단위로 합친 지표가 더 민감할 수 있다.
\item \textbf{학습 쪽 분해능}은 안 쟀다. 유보만 봤다.
\item \textbf{작은 도메인 전용 관문}(노트 529) 여전히 열려 있다.
\end{itemize}

\end{document}
